\chapter{Pancreas}
\index{Pancreas}

\medskip
\noindent\textit{This chapter is a short overview. Persistent or severe abdominal symptoms need clinical assessment.}

\section*{What it is}
The pancreas is an organ behind the stomach. It releases digestive enzymes into the small intestine and produces hormones, including insulin and glucagon, which help control blood glucose.

\section*{Common pancreatic problems}
Pancreatitis is inflammation that may be sudden or long-lasting. Other problems include cysts, blockage of the pancreatic or bile ducts, pancreatic exocrine insufficiency, and cancer. Gallstones, alcohol, some medicines, high triglycerides, inherited conditions, and other illnesses can contribute to pancreatitis.

\section*{Symptoms}
Pancreatic disease may cause upper-abdominal pain that can spread to the back, nausea, vomiting, fever, jaundice, weight loss, greasy or difficult-to-flush stools, or changes in blood glucose. Symptoms overlap with many other conditions, so they do not establish a diagnosis.

\section*{When to Seek Medical Care}
Seek urgent care for sudden severe or persistent abdominal pain, repeated vomiting, fever with worsening pain, jaundice, confusion, fainting, breathing difficulty, or dehydration. New unexplained weight loss, persistent digestive changes, or recurrent pain should be assessed promptly.

\section*{Diagnosis and treatment}
Evaluation may include examination, blood tests, ultrasound, CT, MRI or endoscopic tests. Treatment depends on the cause and may include fluids, pain relief, nutritional support, enzyme replacement, diabetes care, endoscopic or surgical treatment, and cancer care. Avoid alcohol and do not smoke; ask a clinician before taking supplements or medicines marketed for “pancreas cleansing.”

\section*{Sources and further reading}
\begin{itemize}
\item National Institute of Diabetes and Digestive and Kidney Diseases. \textit{Pancreatitis}. \url{https://www.niddk.nih.gov/health-information/digestive-diseases/pancreatitis}
\item National Cancer Institute. \textit{Pancreatic cancer}. \url{https://www.cancer.gov/types/pancreatic}
\end{itemize}
